\chapter{The crisis}
\label{ch:crisis}

\section{A note on writing about the present}

This chapter comes closest to the day it was written, and the closer it comes
the less certain it gets. Cuban statistical publication thinned after 2020;
several series in the \emph{Anuario} stopped appearing or appeared very late;
and the figures for 2025 and 2026 that circulate are in most cases estimates by
independent economists and demographers rather than published data. Where this
chapter gives a recent number it says whose number it is, and the figures for
the last two years should be read as provisional and checked against whatever
has been published since. The direction of every one of them is not in doubt.

\section{A new president, and a new constitution}

Miguel D\'iaz-Canel became president in April 2018 and First Secretary of the
party in April 2021, when Ra\'ul Castro retired. He is the first Cuban leader
since 1959 not named Castro and the first who did not fight in the
insurrection, and neither fact has translated into much independent authority.

His first substantial act was a constitution. The 2019 text, approved in a
February referendum, is a more liberal document than the 1976 one it replaced
and this should be said before the qualification. It recognises private
property and foreign investment as components of the economy rather than as
tolerated exceptions; it restores the presumption of innocence and habeas
corpus; it limits the president to two five-year terms and imposes an age limit
on first election; it re-creates the office of prime minister; it strengthens
the language on non-discrimination. Against that, it retains the Communist
Party as the superior leading force of society and the state, and it retains
the irrevocability of socialism that had been inserted in 2002 to bury the
Varela Project. It is a reform of the state's administration and not of its
politics.

The draft had contained a definition of marriage that would have permitted
same-sex marriage; it was withdrawn after organised opposition from evangelical
churches and referred to a future family code. That code was put to a separate
referendum in September 2022 and passed with about two thirds of the vote,
making Cuba one of the few countries in the world to legalise same-sex marriage
by popular vote. It is a genuine liberal advance and it was delivered by
referendum, which is worth noting in a chapter that is otherwise unrelieved.

\section{The re-tightening}

Between 2017 and 2020 the United States imposed a long series of measures ---
the figure usually cited is 243 --- that reversed the opening and went beyond
it.\unv{} Cruise ship calls were prohibited in June 2019. Individual
people-to-people travel was ended. Remittances were capped and then, in late
2020, effectively cut off when Western Union closed its more than four hundred
Cuban offices after the US Treasury designated the military-linked processor
through which they ran. Title III of Helms--Burton was allowed to come into
force in May 2019 for the first time since 1996, exposing any company using
confiscated property to suit in American courts. And in January 2021, in the
last week of the Trump administration, Cuba was returned to the list of state
sponsors of terrorism, a designation whose practical effect is less in its own
prohibitions than in the behaviour it induces in banks, which de-risk by
refusing Cuban business entirely.

The Biden administration eased parts of this from 2022 --- remittance caps
lifted, group travel restored, consular processing resumed in Havana --- and in
January 2025 removed the terrorism designation, which was reinstated within days
of the change of administration. The net effect over the decade is that the
external environment tightened, that the tightening was concentrated on exactly
the flows that reached Cuban households rather than the Cuban state ---
remittances, independent travel, private accommodation --- and that Cuban
policy toward its own private sector, described in the previous chapter, had
begun tightening first.

Meanwhile the Venezuelan subsidy shrank with Venezuela. Oil deliveries that had
been near 100,000 barrels a day fell to a fraction of that, with the shortfall
covered irregularly by Russia, Mexico and spot purchases the country could not
reliably pay for.

\section{Two years with the borders shut}

Tourism had reached about 4.7 million visitors in 2018. It fell to roughly 1.1
million in 2020 and to something under 400,000 in 2021.\unv{} For an economy in
which tourism was the second or third source of foreign exchange and the first
source of demand for the private sector, closing the borders for two years was
equivalent to closing the private economy. Arrivals recovered to around 2.4
million by 2023 and then stalled below that, at roughly half the pre-pandemic
level, while the Caribbean as a whole recovered fully --- a divergence that has
more to do with blackouts, hotels without water and empty shops than with
American policy.\unv{}

Against this, the pandemic produced the one unambiguous Cuban success of the
period. Cuba's biotechnology sector developed its own COVID-19 vaccines ---
Abdala at the Centre for Genetic Engineering and Biotechnology, Soberana~02 and
Soberana Plus at the Finlay Institute --- ran its own trials, reported
efficacies above 90 per cent, manufactured at national scale, and vaccinated
over 90 per cent of the population with domestically produced product,
including children from the age of two earlier than most of the
world.\unv{} No other country in Latin America did this. Cuba
did it while unable to import syringes reliably.

The qualification belongs with it: the vaccines were not submitted through, or
did not complete, the international prequalification and regulatory pathways
that would have allowed them to be sold at scale abroad, which is the same
failure Chapter~\ref{ch:life} identifies across the whole sector --- excellent
science, no regulatory conversion, and therefore very little money. Cuba
vaccinated itself and could not sell the result.

\section{Tarea Ordenamiento}

On 1 January 2021 --- in the middle of a pandemic, with the borders closed,
tourism at zero and the economy having contracted by about 11 per cent the
previous year --- the government executed the currency unification it had been
preparing since 2011.

The convertible peso was abolished with a six-month conversion window. The
official exchange rate for enterprises went from one peso to the dollar to
twenty-four, a devaluation of 2,300 per cent applied at a stroke to every
balance sheet in the country. Wages and pensions were raised several-fold to
compensate, the minimum wage going to 2,100 pesos. Subsidies were cut and many
regulated prices raised.

The direction was right and had been right for twenty years. Chapter~\ref{ch:dollarreform}
explained why: with two rates twenty-five-fold apart, no Cuban enterprise
account meant anything, and nothing could be reformed because nothing could be
measured. The execution was catastrophic, and the reasons are specific enough to
be instructive.

It was done with no reserves. A fixed rate is a promise to sell dollars at that
rate, and Cuba had no dollars, so the promise was not kept and a parallel rate
appeared immediately.

It was done with no supply response. Wages were multiplied in an economy where
the quantity of goods was falling, which is the definition of an inflation.

It was done with the private sector still repressed --- the licence freeze of
2017 was still in effect, and the small and medium enterprise law was eight
months away --- so the one part of the economy that could have responded to
higher prices with more output was not permitted to.

And it was done at the bottom of the worst external shock since 1993, when the
correct sequence is to stabilise first and unify afterwards.

The results followed. Official consumer price inflation was around 77 per cent
in 2021 and remained in the high double digits for the following two years,
with independent estimates considerably higher for the goods people actually
buy.\unv{} The informal exchange rate, which the reform was meant to eliminate,
detached from the official one and kept going: past 100 pesos to the dollar in
2022, past 200 in 2023, past 300 in 2024, and beyond that since.\unv{} Real
wages and pensions ended lower than before the reform that had multiplied them.
Cuban savings were destroyed for the second time in thirty years.

Chapter~\ref{ch:money} is written against this chapter. The lesson is not that
unification was wrong. It is that a monetary reform without reserves, without a
supply response and without a stabilisation programme in front of it is not a
reform but a devaluation of the population's savings, and that the same state
had demonstrated in 1994--96, with fewer resources and greater competence, how
to do it in the right order.

\section{11 July 2021}

On the morning of Sunday 11 July 2021 people came into the street in San
Antonio de los Ba\~nos, a town of forty thousand south-west of Havana. It was
filmed and put on Facebook. Within hours there were demonstrations in dozens of
towns and cities across the island --- the counts run from about forty to over
sixty locations --- including Havana, Santiago, Camag\"uey, Holgu\'in,
Cienfuegos, Matanzas and Palma Soriano.\unv{}

Nothing like it had happened since 1959. The \emph{maleconazo} of 1994 was one
city on one afternoon. This was national, simultaneous, and organised by nobody
--- which is exactly what made it uncontainable and, for the government,
incomprehensible. The immediate causes were blackouts in a July heatwave, empty
pharmacies during the worst COVID wave Cuba had, food queues, and the price
shock of the ordenamiento. The demonstrators were overwhelmingly young, poor,
and disproportionately from the Black and mixed-race neighbourhoods that
Chapter~\ref{ch:dollarreform} identified as the households without remittances
--- which is to say, from the constituency the revolution has always counted as
its own.

The response was mobile internet shut off nationally within hours; the
president's televised instruction, the \emph{orden de combate}, calling
revolutionaries into the street; deployment of special forces and of civilian
brigades; and mass arrests. One protester, Diubis Laurencio Tejeda, was shot
dead in La G\"uinera in Havana. Cubalex and Justicia~11J, the two monitoring
groups that did the counting because nobody official would \citep{justicia11j}, documented 1,481
people deprived of liberty in connection with the protests, of whom 57 were
under eighteen. Around 600 were prosecuted; at least 141 were charged with
sedition; and sentences ran from four years to thirty. Three years on, several
hundred were still in prison. The prosecutions were not of an organisation,
because there was no organisation; they were of individuals who had walked in
the street.

This is the central political fact of the period and Part III has to be answerable
to it. A government whose claim to legitimacy rests on having lifted the poor
and the Black out of the republic's neglect responded to a demonstration by the
poor and the Black with twenty-year sentences. Whatever the design in Part III
proposes about reconciliation, amnesty and the sequencing of political opening,
11 July 2021 and its prisoners are the case it has to be able to handle.

\section{The state's own retreat}

In August 2021, six weeks after the protests, the government did the thing it
had refused to do for a decade: Decree-Laws 46 to 49 created small and medium
enterprises with legal personality. A Cuban private business could at last be a
firm --- own assets, sign contracts, hire, and import and export through a state
intermediary. Some eleven or twelve thousand were approved over the following
three years, the large majority private, and they rapidly became a significant
share of the country's imports of food and consumer goods.\unv{}

The reform was real and it was also, transparently, a surrender: the state
legalised private importers because it could no longer stock its own shops.
The same logic produced the partial re-dollarisation of retail from 2024 ---
new state shops selling in dollars only, to capture remittances --- which
reintroduced, three years after abolishing it, a two-currency retail economy
with the currency of the poor on the wrong side of it.

And the enterprises arrived into a policy environment that turned on them
almost at once. From 2024 came price caps on private traders, restrictions on
what could be imported, higher taxes, and a public argument from the leadership
that the new firms were the cause of inflation rather than a response to it. The
pattern of 1996 and of 2017 repeated on a two-year cycle.

\section{The lights}

Cuba's electricity system reached the end of its life in 2024.

The generating fleet is a small number of large oil-fired thermal units, most
built between the 1960s and the 1980s, running past their design lives on heavy
domestic crude they were not built for, without maintenance, without spares, and
frequently without fuel. Through 2023 and 2024 daily deficits of a third or more
of demand were routine, and blackouts of twelve to twenty hours were normal
outside Havana and increasingly normal inside it.

The scale of the retreat is visible in the output. Cuba generated on the order
of 20 terawatt-hours in 2018; in 2023 it generated 15.3, of which 3.7 --- close
to a quarter --- was lost in transmission and distribution before reaching
anybody.\unv{} Consumption per head fell from about 1,856 kilowatt-hours in
2018 to 1,387 in 2023.

On 18 October 2024 the Antonio Guiteras plant, the largest single unit in the
country, tripped, and the entire national grid collapsed. Cuba went dark ---
eleven million people, every province, no exceptions. Restoration took several
days and failed twice during the attempt. It happened again in November after
Hurricane Rafael, and again in December, and again in 2025.\unv{} A national
blackout is not a shortage; it is the failure of the system as a system, and
once a grid has collapsed it collapses more easily.

The government's answer is solar, financed and built by China: a programme
announced around 2024 of dozens of photovoltaic parks with a first target near a
thousand megawatts, of which several hundred megawatts were reported energised
during 2025.\unv{} The direction is right --- Chapter~\ref{ch:energy} argues it
is the only direction available --- but photovoltaic capacity without storage
does not address a demand peak that occurs after sunset, and it does not by
itself restart a thermal fleet that must still cover the night.

Everything else in this chapter is downstream of the lights. Without stable
power there is no water pumping, no refrigeration, no food processing, no
tourism worth the name, no hospital that can plan a list, no remote work, and no
industry. Chapter~\ref{ch:premises} therefore treats electricity as the binding
constraint on the whole design and puts it before every other sector.

\section{The exodus}

The largest emigration in Cuban history took place between 2021 and 2025, and it
dwarfs the ones this book has already described.

The route ran through Nicaragua, which removed visa requirements for Cubans in
November 2021, and then overland to the United States border. In the American
fiscal year 2022 US authorities recorded something over 220,000 encounters with
Cuban nationals, and about 200,000 the following year; a humanitarian parole
programme opened in January 2023 admitted tens of thousands more by air; and
substantial numbers went to Spain under a 2022 descent-based nationality law, to
Mexico, Uruguay, Serbia and elsewhere.\unv{} Reasonable estimates of total
departures over 2021--2025 run from about one million to something approaching
two, against a population that had been about 11.2 million.

The official statistics office put the resident population at 9.7 million in
2024 --- a fall of 1.4 million in four years, and the first time it has been
below ten million since the 1980s --- with the 2024 census postponed.
Independent Cuban demographers, working from birth, death and border data, put
the true resident figure lower still, one study at around eight million, and
Albizu-Campos estimates roughly 1.79 million departures over 2021--24
alone \citep{onei}.\unv{} The uncertainty is large and the fact is not: Cuba has lost
something like a tenth to a fifth of its people in five years, and they were
disproportionately of working and child-bearing age.

This lands on a demography that was already among the oldest in the hemisphere.
Fertility has been below replacement since 1978 --- there is no cohort in the
pipeline --- and by the mid-2020s the total fertility rate was around 1.2 to
1.4, among the lowest in the world, with more than a quarter of the population
over sixty.\unv{} A pension system, a health system and a tax base all depend on
a ratio of workers to dependants, and that ratio has now moved further and
faster in Cuba than in any country not at war.

Chapter~\ref{ch:welfare} treats this as the hardest number in the book, and
Chapter~\ref{ch:talent} draws the conclusion that follows from it: a country
that has lost this many working-age people cannot be designed as though labour
were abundant. Every proposal in Part III is labour-constrained.

\section{Where the country stands}

In 2026 Cuba is poorer, older, smaller, darker and hungrier than at any point
since 1994, and by several measures than at any point since 1959.

Output has contracted in most years since 2019 and the cumulative fall is
large.\unv{} Sugar, which was 8.5 million tonnes in 1970 and 4 million as
recently as the 2000s, produced a harvest measured in the low hundreds of
thousands of tonnes in the mid-2020s --- less than in the nineteenth century ---
and Cuba now imports sugar.\unv{} The rationed basket is delivered late and
incomplete, and in 2023 the government asked the World Food Programme for
assistance with powdered milk for children, which it had never previously
done.\unv{} Medicines are absent from pharmacies at rates the health ministry
itself has acknowledged. Doctors and nurses are emigrating, and so are teachers.
The blackouts continue.

Against that: the schools are open, the clinics are staffed however thinly, the
vaccination programme still runs, violent crime remains low by regional
standards, and the state has not lost control of the territory. The floor that
held in 1993 is thinner and it is still there. It is the thing Part III is
written to save.

\begin{ledger}{the crisis, 2018--2026}
\emph{Achieved.} A constitution that recognises private property and restores
habeas corpus, and a family code legalising same-sex marriage passed by
referendum. Three domestically developed COVID-19 vaccines and a population
vaccinated with them, alone in Latin America. Private firms with legal
personality at last, after a decade of refusal. A currency unification that,
whatever its execution, ended a twenty-six-year accounting fiction. Solar
capacity begun. And a state that has kept schools open and clinics staffed
through a contraction it has not been able to stop.

\emph{Cost.} A monetary reform executed with no reserves, no supply response and
the private sector still shut, which destroyed the population's savings for the
second time in thirty years. Twelve-to-twenty-hour blackouts and four national
grid collapses. Sugar production below nineteenth-century levels and an appeal
to the World Food Programme for children's milk. The largest emigration in
Cuban history --- between one and two million people in five years, from a
population of eleven --- onto the oldest demography in the hemisphere. And, on
11 July 2021, sentences of up to thirty years, including of minors,
handed to poor and mostly Black Cubans for walking in the street to say there
was no food and no electricity.
\end{ledger}
